\documentclass[letterpaper,12pt]{article}

\usepackage{fancyhdr}
\usepackage[most]{tcolorbox}
\usepackage{longtable}
\usepackage{makecell}
\usepackage{tikz}
\usepackage{geometry}
\usepackage{array}
\usepackage{multicol}
\usepackage[table]{xcolor}
\usepackage{polyglossia}
\setmainlanguage{english}
\setotherlanguage{arabic}
\setotherlanguage{hebrew}
\usepackage{fontspec}
\setmainfont{Charis SIL}
\newfontfamily\arabicfont[Script=Arabic,Scale=1.2]{Amiri}
\newfontfamily\hebrewfont{Ezra SIL}[Script=Hebrew]
\geometry{margin=1.5cm}
\usepackage{booktabs}
\usepackage{graphicx}
\usepackage{multirow}

\setlength{\headheight}{14.49998pt}
\addtolength{\topmargin}{-2.49998pt}
\pagestyle{fancy}
\fancyhf{}
\fancyfoot[C]{\thepage}
\fancyhead[R]{\textit{Why Am I Forced to Learn a Useless Language?}}
\renewcommand{\headrulewidth}{0pt}

\usetikzlibrary{shapes,arrows,decorations.pathmorphing}

% Custom colors
\definecolor{headercolor}{RGB}{70,130,180}
\definecolor{boxcolor}{RGB}{240,248,255}
\definecolor{accentcolor}{RGB}{220,20,60}
\definecolor{tableheader}{RGB}{220,220,220}
\definecolor{dialectcolor}{RGB}{34,139,34}

\begin{document}

\title{\textbf{\Large Arabic Phrase Analyser}\\
\large Why Am I Forced to Learn a Useless Language?\\
\normalsize \textarabic{لم اقحم في مخي لغة لا طائلة منها في الدنيا}}
\author{Igor Deruga}
\date{}
\maketitle

% ======================== Phrase Display ========================
\begin{tcolorbox}[colback=boxcolor,colframe=headercolor,title=\textbf{Arabic Phrase with Full Diacritics},breakable]
\centering
\large\textarabic{لِمَ أُقْحَمُ فِي مُخِّي لُغَةً لَا طَائِلَةَ مِنْهَا فِي الدُّنْيَا؟}
\end{tcolorbox}

% ======================== Translations ========================
\section{English Translation}
\begin{tcolorbox}[colback=white,colframe=accentcolor,breakable]
\textbf{Literal:} Why am-I-forced in my-brain a-language no benefit from-it in the-world?\\
\textit{[Arabic order retained for direct mapping]}\\[0.5em]
\textbf{Adapted:} Why am I forced to cram into my brain a language from which there is no benefit in this world?
\end{tcolorbox}

% ======================== Detailed Word Analysis ========================
\section{Detailed Word Analysis}

\subsection{\textarabic{لِمَ} — [lima]}
\begin{tabular}{p{3cm}p{10cm}}
\toprule
\textbf{Translation} & why? / for what? \\
\textbf{Root} & Interrogative particle, not derived from root \\
\textbf{Pattern} & Composite: \textarabic{لِ} (li-) + \textarabic{مَا} (maË) \\
\textbf{Grammar} & Interrogative particle asking for reason/cause; \textarabic{لِ} (preposition ``for'') + \textarabic{مَا} (``what''), with elided final alif \\
\midrule
\textbf{Examples} & \makecell[l]{\parbox{9.5cm}{
1. \textarabic{لِمَ تَبْكِي؟} — Why are you crying? [lima tabkiË]\\
2. \textarabic{لِمَ لَمْ تَأْتِ أَمْسِ؟} — Why didn't you come yesterday? [lima lam taÊ"ti Ê"amsi]\\
3. \textarabic{لِمَ يُعَامِلُنَا هَكَذَا؟} — Why does he treat us this way? [lima juÊ•aËmilnunaË haËkaðaË]
}} \\
\midrule
\textbf{Synonyms} & \textarabic{لِمَاذَا} (limaËðaË — why), \textarabic{لِأَيِّ شَيْءٍ} (liÊ"ajji ʃajÊ"in — for what thing) \\
\textbf{Etymology} & Combination of preposition \textarabic{لِ} (for) + interrogative \textarabic{مَا} (what) \\
\bottomrule
\end{tabular}

\subsubsection*{Usage Notes}
\begin{itemize}
\item \textarabic{لِمَ} is a contraction of \textarabic{لِمَاذَا}, more literary and concise
\item Often used in rhetorical questions expressing frustration or complaint
\item The \textarabic{لِ} prefix adds purposive meaning (``for what purpose'')
\item In Modern Standard Arabic, \textarabic{لِمَاذَا} is more common in speech
\end{itemize}

\subsection{\textarabic{أُقْحَمُ} — [ʔuqħamu]}
\begin{tabular}{p{3cm}p{10cm}}
\toprule
\textbf{Translation} & I am forced / I am thrust / I am crammed \\
\textbf{Root} & \textarabic{ق-ح-م} (q-ħ-m) \\
\textbf{Pattern} & \textarabic{أُفْعَلُ} (Ê"ufÊ•alu) — Form IV passive imperfect \\
\textbf{Grammar} & Present passive verb, 1st person singular; indicates being forced into something unwillingly \\
\midrule
\textbf{Examples} & \makecell[l]{\parbox{9.5cm}{
1. \textarabic{أَقْحَمَ نَفْسَهُ فِي الْمَشَاكِلِ} — He thrust himself into problems [Ê"aqħama nafsahu fiË l-maʃaËkili]\\
2. \textarabic{لَا تُقْحِمْنِي فِي هَذَا الْأَمْرِ} — Don't drag me into this matter [laË tuqħimniË fiË haËða l-Ê"amri]\\
3. \textarabic{أُقْحِمَ الطَّالِبُ فِي امْتِحَانٍ صَعْبٍ} — The student was thrust into a difficult exam [Ê"uqħima tˤ-tˤaËlibu fiË mtiħaËnin sˤaÊ•bin]
}} \\
\midrule
\textbf{Synonyms} & \textarabic{أُجْبَرُ} (forced), \textarabic{أُرْغَمُ} (compelled), \textarabic{أُحْشَرُ} (stuffed/crammed) \\
\textbf{Etymology} & Root meaning: to thrust oneself, to plunge, to force entry \\
\bottomrule
\end{tabular}

\subsection{Conjugation — \textarabic{أَقْحَمَ} (Form IV)}
\par{\large\textbf{Verbal Noun}: \textarabic{إِقْحَام} /Ê"iqħaËm/ — forcing, thrusting}

\begin{longtable}{|>{\raggedright}p{3.5cm}|p{5cm}|p{5cm}|}
\hline
\textbf{Person} & \textbf{Perfect (Past)} & \textbf{Imperfect (Present)} \\
\hline
\endfirsthead
\hline
\textbf{Person} & \textbf{Perfect (Past)} & \textbf{Imperfect (Present)} \\
\hline
\endhead
\textbf{3rd m. sg.} & \textarabic{أَقْحَمَ} /Ê"aqħama/ & \textarabic{يُقْحِمُ} /juqħimu/ \\
\hline
\textbf{3rd f. sg.} & \textarabic{أَقْحَمَتْ} /Ê"aqħamat/ & \textarabic{تُقْحِمُ} /tuqħimu/ \\
\hline
\textbf{3rd m. dual} & \textarabic{أَقْحَمَا} /Ê"aqħamaË/ & \textarabic{يُقْحِمَانِ} /juqħimaËni/ \\
\hline
\textbf{3rd f. dual} & \textarabic{أَقْحَمَتَا} /Ê"aqħamataË/ & \textarabic{تُقْحِمَانِ} /tuqħimaËni/ \\
\hline
\textbf{3rd m. pl.} & \textarabic{أَقْحَمُوا} /Ê"aqħamuË/ & \textarabic{يُقْحِمُونَ} /juqħimuËna/ \\
\hline
\textbf{3rd f. pl.} & \textarabic{أَقْحَمْنَ} /Ê"aqħamna/ & \textarabic{يُقْحِمْنَ} /juqħimna/ \\
\hline
\textbf{2nd m. sg.} & \textarabic{أَقْحَمْتَ} /Ê"aqħamta/ & \textarabic{تُقْحِمُ} /tuqħimu/ \\
\hline
\textbf{2nd f. sg.} & \textarabic{أَقْحَمْتِ} /Ê"aqħamti/ & \textarabic{تُقْحِمِينَ} /tuqħimiËna/ \\
\hline
\textbf{2nd dual (m./f.)} & \textarabic{أَقْحَمْتُمَا} /Ê"aqħamtumaË/ & \textarabic{تُقْحِمَانِ} /tuqħimaËni/ \\
\hline
\textbf{2nd m. pl.} & \textarabic{أَقْحَمْتُمْ} /Ê"aqħamtum/ & \textarabic{تُقْحِمُونَ} /tuqħimuËna/ \\
\hline
\textbf{2nd f. pl.} & \textarabic{أَقْحَمْتُنَّ} /Ê"aqħamtunna/ & \textarabic{تُقْحِمْنَ} /tuqħimna/ \\
\hline
\textbf{1st sg.} & \textarabic{أَقْحَمْتُ} /Ê"aqħamtu/ & \textarabic{أُقْحِمُ} /Ê"uqħimu/ \\
\hline
\textbf{1st pl.} & \textarabic{أَقْحَمْنَا} /Ê"aqħamnaË/ & \textarabic{نُقْحِمُ} /nuqħimu/ \\
\hline
\end{longtable}

\subsubsection*{Passive Voice}
\begin{longtable}{|>{\raggedright}p{3.5cm}|p{5cm}|p{5cm}|}
\hline
\textbf{Person} & \textbf{Perfect (Past)} & \textbf{Imperfect (Present)} \\
\hline
\textbf{3rd m. sg.} & \textarabic{أُقْحِمَ} /Ê"uqħima/ & \textarabic{يُقْحَمُ} /juqħamu/ \\
\hline
\textbf{1st sg.} & \textarabic{أُقْحِمْتُ} /Ê"uqħimtu/ & \textarabic{أُقْحَمُ} /Ê"uqħamu/ \\
\hline
\textbf{1st pl.} & \textarabic{أُقْحِمْنَا} /Ê"uqħimnaË/ & \textarabic{نُقْحَمُ} /nuqħamu/ \\
\hline
\end{longtable}

\subsection{\textarabic{فِي} — [fiË]}
\begin{tabular}{p{3cm}p{10cm}}
\toprule
\textbf{Translation} & in / into / within \\
\textbf{Root} & Preposition, not root-based \\
\textbf{Pattern} & Simple preposition \\
\textbf{Grammar} & Preposition governing genitive case; indicates location, containment, or involvement \\
\midrule
\textbf{Examples} & \makecell[l]{\parbox{9.5cm}{
1. \textarabic{فِي الْبَيْتِ} — In the house [fiË l-bajti]\\
2. \textarabic{يَعْمَلُ فِي الشَّرِكَةِ} — He works in the company [jaÊ•malu fiË Êƒ-ʃarikati]\\
3. \textarabic{فِي الْمَسَاءِ} — In the evening [fiË l-masaËÊ"i]
}} \\
\midrule
\textbf{Synonyms} & \textarabic{بِ} (with/in), \textarabic{دَاخِلَ} (inside), \textarabic{ضِمْنَ} (within) \\
\textbf{Etymology} & Ancient Arabic preposition, cognate with Hebrew \texthebrew{בְּ} (be-) \\
\bottomrule
\end{tabular}

\subsection{\textarabic{مُخِّي} — [muxxiË]}
\begin{tabular}{p{3cm}p{10cm}}
\toprule
\textbf{Translation} & my brain \\
\textbf{Root} & \textarabic{م-خ-خ} (m-x-x) \\
\textbf{Pattern} & \textarabic{فُعّ} (fuÊ•Ê•) with possessive suffix \\
\textbf{Grammar} & Noun, masculine, genitive case (after preposition \textarabic{فِي}), with 1st person singular possessive suffix \textarabic{ـِي} \\
\midrule
\textbf{Examples} & \makecell[l]{\parbox{9.5cm}{
1. \textarabic{الْمُخُّ الْبَشَرِيُّ} — The human brain [al-muxxu l-baʃarijju]\\
2. \textarabic{يُؤْلِمُنِي مُخِّي} — My brain hurts me [juÊ"limniË muxxiË]\\
3. \textarabic{مُخُّ الْإِنْسَانِ مُعَقَّدٌ} — The human brain is complex [muxxu l-Ê"insaËni muÊ•aqqadun]
}} \\
\midrule
\textbf{Synonyms} & \textarabic{دِمَاغ} (dimaËÉ£ — brain), \textarabic{عَقْل} (Ê•aql — mind/intellect) \\
\textbf{Etymology} & Root \textarabic{م-خ-خ} refers to marrow, brain substance \\
\bottomrule
\end{tabular}

\subsubsection*{Full Declension with Possessive Suffixes}
\begin{tabular}{|c|c|c|c|}
\hline
\textbf{Case} & \textbf{Without Suffix} & \textbf{My brain} & \textbf{Your (m.) brain} \\
\hline
Nominative & \textarabic{مُخٌّ} /muxxun/ & \textarabic{مُخِّي} /muxxiË/ & \textarabic{مُخُّكَ} /muxxuka/ \\
\hline
Accusative & \textarabic{مُخًّا} /muxxan/ & \textarabic{مُخِّي} /muxxiË/ & \textarabic{مُخَّكَ} /muxxaka/ \\
\hline
Genitive & \textarabic{مُخٍّ} /muxxin/ & \textarabic{مُخِّي} /muxxiË/ & \textarabic{مُخِّكَ} /muxxika/ \\
\hline
\end{tabular}

\subsection{\textarabic{لُغَةً} — [luÉ£atan]}
\begin{tabular}{p{3cm}p{10cm}}
\toprule
\textbf{Translation} & a language \\
\textbf{Root} & \textarabic{ل-غ-و} (l-É£-w) \\
\textbf{Pattern} & \textarabic{فُعَلَة} (fuÊ•ala) \\
\textbf{Grammar} & Noun, feminine, accusative case (object of passive verb), indefinite with tanwin \\
\midrule
\textbf{Examples} & \makecell[l]{\parbox{9.5cm}{
1. \textarabic{أَتَعَلَّمُ لُغَةً جَدِيدَةً} — I am learning a new language [Ê"ataÊ•allamu luÉ£atan dÊ'adiËdatan]\\
2. \textarabic{اللُّغَةُ الْعَرَبِيَّةُ} — The Arabic language [al-luÉ£atu l-Ê•arabijjatu]\\
3. \textarabic{يَتَكَلَّمُ ثَلَاثَ لُغَاتٍ} — He speaks three languages [jatakalamu θalaËθa luÉ£aËtin]
}} \\
\midrule
\textbf{Synonyms} & \textarabic{لِسَان} (lisaËn — tongue/language), \textarabic{كَلَام} (kalaËm — speech) \\
\textbf{Etymology} & From root meaning ``to speak, to use language''; related to Hebrew \texthebrew{לָשׁוֹן} (lashon) \\
\bottomrule
\end{tabular}

\subsubsection*{Full Declension}
\begin{tabular}{|c|c|c|c|}
\hline
\textbf{Number/Case} & \textbf{Indefinite} & \textbf{Definite} & \textbf{Plural Indefinite} \\
\hline
Nominative sg. & \textarabic{لُغَةٌ} /luÉ£atun/ & \textarabic{اللُّغَةُ} /al-luÉ£atu/ & \textarabic{لُغَاتٌ} /luÉ£aËtun/ \\
\hline
Accusative sg. & \textarabic{لُغَةً} /luÉ£atan/ & \textarabic{اللُّغَةَ} /al-luÉ£ata/ & \textarabic{لُغَاتٍ} /luÉ£aËtin/ \\
\hline
Genitive sg. & \textarabic{لُغَةٍ} /luÉ£atin/ & \textarabic{اللُّغَةِ} /al-luÉ£ati/ & \textarabic{لُغَاتٍ} /luÉ£aËtin/ \\
\hline
\end{tabular}

\subsection{\textarabic{لَا} — [laË]}
\begin{tabular}{p{3cm}p{10cm}}
\toprule
\textbf{Translation} & no / not \\
\textbf{Root} & Negative particle, not root-based \\
\textbf{Pattern} & Simple particle \\
\textbf{Grammar} & Negative particle; when followed by indefinite noun in accusative, acts as \textarabic{لَا النَّافِيَة لِلْجِنْسِ} (categorical negation) \\
\midrule
\textbf{Examples} & \makecell[l]{\parbox{9.5cm}{
1. \textarabic{لَا أَحَدَ هُنَا} — No one is here [laË Ê"aħada hunaË]\\
2. \textarabic{لَا شَكَّ فِي ذَلِكَ} — No doubt about that [laË ʃakka fiË Ã°aËlika]\\
3. \textarabic{لَا فَائِدَةَ مِنْ هَذَا} — No benefit from this [laË faËÊ"idata min haËðaË]
}} \\
\midrule
\textbf{Synonyms} & \textarabic{لَيْسَ} (negation verb), \textarabic{مَا} (negative particle in some contexts) \\
\textbf{Etymology} & Ancient Semitic negative particle \\
\bottomrule
\end{tabular}

\subsection{\textarabic{طَائِلَةَ} — [tˤaËÊ"ilata]}
\begin{tabular}{p{3cm}p{10cm}}
\toprule
\textbf{Translation} & benefit / use / advantage \\
\textbf{Root} & \textarabic{Ø·-ÙˆÙŽ-ل} (tˤ-w-l) \\
\textbf{Pattern} & \textarabic{فَاعِلَة} (faËÊ•ila) — Form I active participle, feminine \\
\textbf{Grammar} & Noun/participle, feminine, accusative case (predicate after \textarabic{لَا}), indefinite \\
\midrule
\textbf{Examples} & \makecell[l]{\parbox{9.5cm}{
1. \textarabic{لَا طَائِلَةَ مِنْ هَذَا الْعَمَلِ} — No benefit from this work [laË tˤaËÊ"ilata min haËða l-Ê•amali]\\
2. \textarabic{طَائِلٌ تَحْتَ يَدِهِ} — Abundant wealth in his possession [tˤaËÊ"ilun taħta jadihi]\\
3. \textarabic{مَالٌ طَائِلٌ} — Abundant/considerable wealth [maËlun tˤaËÊ"ilun]
}} \\
\midrule
\textbf{Synonyms} & \textarabic{فَائِدَة} (benefit), \textarabic{نَفْع} (utility), \textarabic{جَدْوَى} (usefulness) \\
\textbf{Etymology} & From root meaning ``to reach, to attain, to be abundant'' \\
\bottomrule
\end{tabular}

\subsubsection*{Usage Notes}
\begin{itemize}
\item \textarabic{طَائِلَة} primarily appears in the expression \textarabic{لَا طَائِلَةَ مِنْ} (no benefit from)
\item Can mean both ``benefit'' and ``abundance/wealth'' depending on context
\item Literary and formal register; rarely used in colloquial speech
\item The active participle form suggests ``reaching/attaining'' quality
\end{itemize}

\subsection{\textarabic{مِنْهَا} — [minhaË]}
\begin{tabular}{p{3cm}p{10cm}}
\toprule
\textbf{Translation} & from it / from her \\
\textbf{Root} & Preposition + pronoun \\
\textbf{Pattern} & \textarabic{مِنْ} + \textarabic{ـهَا} \\
\textbf{Grammar} & Preposition \textarabic{مِنْ} (from) + 3rd person feminine singular pronoun suffix \textarabic{ـهَا} referring to \textarabic{لُغَة} \\
\midrule
\textbf{Examples} & \makecell[l]{\parbox{9.5cm}{
1. \textarabic{أَخَذْتُ مِنْهَا كِتَابًا} — I took a book from her/it [Ê"axaðtu minhaË kitaËban]\\
2. \textarabic{تَعَلَّمْتُ مِنْهَا الْكَثِيرَ} — I learned much from her/it [taÊ•allamtu minhaË l-kaθiËra]\\
3. \textarabic{لَا بُدَّ مِنْهَا} — It is inevitable/necessary [laË budda minhaË]
}} \\
\midrule
\textbf{Related Forms} & \textarabic{مِنْهُ} (from him/it m.), \textarabic{مِنْكَ} (from you m.), \textarabic{مِنِّي} (from me) \\
\textbf{Etymology} & \textarabic{مِنْ} is ancient preposition indicating source or origin \\
\bottomrule
\end{tabular}

\subsection{\textarabic{فِي الدُّنْيَا} — [fiË d-dunjaË]}
\begin{tabular}{p{3cm}p{10cm}}
\toprule
\textbf{Translation} & in the world / in this life \\
\textbf{Root} & \textarabic{د-ن-و} (d-n-w) \\
\textbf{Pattern} & \textarabic{فُعْلَى} (funlaË) — feminine comparative/superlative form \\
\textbf{Grammar} & Feminine noun, genitive case (after preposition), definite with article \textarabic{الـ} \\
\midrule
\textbf{Examples} & \makecell[l]{\parbox{9.5cm}{
1. \textarabic{الدُّنْيَا وَالْآخِرَة} — This world and the hereafter [ad-dunjaË wa-l-Ê"aËxira]\\
2. \textarabic{أَمُورُ الدُّنْيَا} — Worldly matters [Ê"umuËru d-dunjaË]\\
3. \textarabic{حَيَاةُ الدُّنْيَا قَصِيرَةٌ} — The life of this world is short [ħajaËtu d-dunjaË qasˤiËratun]
}} \\
\midrule
\textbf{Synonyms} & \textarabic{الْعَالَم} (al-Ê•aËlam — the world), \textarabic{الْحَيَاة} (al-ħajaËt — life) \\
\textbf{Etymology} & From root meaning ``to be near, low''; \textarabic{الدُّنْيَا} means ``the nearer/lower [life]'' as opposed to the afterlife \\
\bottomrule
\end{tabular}

\subsubsection*{Full Declension}
\begin{tabular}{|c|c|c|}
\hline
\textbf{Case} & \textbf{With Article} & \textbf{Meaning} \\
\hline
Nominative & \textarabic{الدُّنْيَا} /ad-dunjaË/ & the world (subject) \\
\hline
Accusative & \textarabic{الدُّنْيَا} /ad-dunjaË/ & the world (object) \\
\hline
Genitive & \textarabic{الدُّنْيَا} /ad-dunjaË/ & the world (after prep.) \\
\hline
\end{tabular}

\textit{Note: \textarabic{الدُّنْيَا} is a diptote (partially declinable), appearing the same in all three cases}

% ======================== Phrase Analysis ========================
\section{Phrase Analysis}
\begin{tcolorbox}[colback=boxcolor,colframe=headercolor,breakable]
\textbf{Grammatical Structure:}\\
Interrogative particle + passive present verb (1st person) + preposition + possessed noun (genitive) + indefinite noun (accusative) + negative particle + feminine noun/participle (accusative) + prepositional phrase with pronoun + prepositional phrase with definite noun \\[0.5em]

\textbf{Key Grammar Points:}
\begin{itemize}
\item \textbf{Rhetorical Question Structure:} \textarabic{لِمَ} opens a complaint-style rhetorical question expressing frustration
\item \textbf{Passive Voice:} \textarabic{أُقْحَمُ} (Form IV passive) emphasizes the speaker's lack of agency — being forced unwillingly
\item \textbf{Metaphorical Expression:} \textarabic{فِي مُخِّي} (in my brain) is a vivid metaphor for forced memorization/learning
\item \textbf{\textarabic{لَا} Construction:} The negative particle with \textarabic{طَائِلَةَ} in accusative creates categorical negation (no benefit whatsoever)
\item \textbf{Relative Clause:} \textarabic{لَا طَائِلَةَ مِنْهَا} functions as an adjectival phrase modifying \textarabic{لُغَةً}
\item \textbf{Emphatic Context:} \textarabic{فِي الدُّنْيَا} adds emphasis — not just "no benefit" but "no benefit in this entire world"
\item \textbf{Register:} The phrase mixes formal classical structures (\textarabic{أُقْحَمُ}, \textarabic{طَائِلَةَ}) with emotional complaint
\item \textbf{Pragmatic Function:} Expresses strong objection to mandatory language learning perceived as useless
\end{itemize}
\end{tcolorbox}

\subsection{Syntactic Breakdown}
\begin{tcolorbox}[colback=white,colframe=headercolor,breakable]
\begin{enumerate}
\item \textbf{\textarabic{لِمَ}} — Interrogative adverb (asking for reason/purpose)
\item \textbf{\textarabic{أُقْحَمُ}} — Main verb: passive imperfect, 1st person singular
\item \textbf{\textarabic{فِي مُخِّي}} — Prepositional phrase indicating location/manner
  \begin{itemize}
  \item \textarabic{فِي} — Preposition
  \item \textarabic{مُخِّي} — Noun (genitive after preposition) + 1st person possessive suffix
  \end{itemize}
\item \textbf{\textarabic{لُغَةً}} — Direct object (accusative indefinite)
\item \textbf{\textarabic{لَا طَائِلَةَ مِنْهَا}} — Adjectival relative clause modifying \textarabic{لُغَةً}
  \begin{itemize}
  \item \textarabic{لَا} — Categorical negative particle
  \item \textarabic{طَائِلَةَ} — Predicate noun (accusative after \textarabic{لَا})
  \item \textarabic{مِنْهَا} — Prepositional phrase with pronoun referring back to \textarabic{لُغَةً}
  \end{itemize}
\item \textbf{\textarabic{فِي الدُّنْيَا}} — Adverbial prepositional phrase (scope/domain)
\end{enumerate}
\end{tcolorbox}

% ======================== Similar Phrases for Practice ========================
\section{Similar Phrases for Practice}

\begin{enumerate}
\item \textarabic{لِمَ أُجْبَرُ عَلَى دِرَاسَةِ مَادَّةٍ لَا فَائِدَةَ مِنْهَا فِي حَيَاتِي؟}\\
Why am I forced to study a subject from which there is no benefit in my life?\\
{[lima Ê"udÊ'baru Ê•alaË diraËsati maËddatin laË faËÊ"idata minhaË fiË ħajaËtiË]}

\item \textarabic{لِمَ نُكَلَّفُ بِحِفْظِ مَعْلُومَاتٍ لَا نَفْعَ مِنْهَا فِي الْعَمَلِ؟}\\
Why are we tasked with memorizing information from which there is no use in work?\\
{[lima nukallafu biħifðˤi maÊ•luËmaËtin laË nafÊ•a minhaË fiË l-Ê•amali]}

\item \textarabic{لِمَ يُفْرَضُ عَلَيْنَا تَعَلُّمُ شَيْءٍ لَا جَدْوَى مِنْهُ فِي الْمُسْتَقْبَلِ؟}\\
Why is learning something from which there is no usefulness in the future imposed on us?\\
{[lima jufradˤu Ê•alajnaË taÊ•allumu ʃajÊ"in laË dÊ'adwaË minhu fiË l-mustaqbali]}

\item \textarabic{لِمَ أُلْزَمُ بِقِرَاءَةِ كُتُبٍ لَا طَائِلَةَ تَحْتَهَا لِلطُّلَّابِ؟}\\
Why am I obliged to read books under which there is no benefit for students?\\
{[lima Ê"ulzamu biqiraËÊ"ati kutubin laË tˤaËÊ"ilata taħtahaË li-tˤ-tˤullaËbi]}

\item \textarabic{لِمَاذَا نُضْطَرُّ إِلَى حِفْظِ نُصُوصٍ لَا قِيمَةَ لَهَا فِي الْوَاقِعِ؟}\\
Why must we be compelled to memorize texts that have no value in reality?\\
{[limaËðaË nudˤtˤarru Ê"ilaË ħifðˤi nusˤuËsˤin laË qiËmata lahaË fiË l-waËqiÊ•i]}
\end{enumerate}

% ======================== Levantine Dialect Version ========================
\section{Levantine (Shaami) Arabic Dialect}

\begin{tcolorbox}[colback=white,colframe=dialectcolor,title=\textbf{Levantine Version},breakable]
\textarabic{لَيْشْ عَمْ يِحْشُرُوا بِمُخِّي لُغَة مَا إِلهَا فَايْدَة بِالدُّنْيِي؟}\\[0.3em]
\textbf{Phonetic:} [leËʃ Ê•am jiħʃuruË bmuxxiË luÉ£a maË Ê"ilha faËjde bi-d-diniË]\\[0.3em]
\textbf{Translation:} Why are they cramming into my brain a language that has no benefit in this world?
\end{tcolorbox}

\textbf{Key Dialectal Changes:}
\begin{itemize}
\item \textarabic{لِمَ} â†' \textarabic{لَيْشْ} (leËʃ) — common Levantine interrogative for "why"
\item \textarabic{أُقْحَمُ} â†' \textarabic{عَمْ يِحْشُرُوا} (Ê•am jiħʃuruË) — active voice "they are cramming" with continuous marker \textarabic{عَمْ}
\item Passive construction replaced with active 3rd person plural (impersonal "they")
\item \textarabic{فِي مُخِّي} â†' \textarabic{بِمُخِّي} (bmuxxiË) — \textarabic{بِ} replaces \textarabic{فِي} in dialect
\item \textarabic{لُغَةً} â†' \textarabic{لُغَة} (luÉ£a) — no case endings in dialect
\item \textarabic{لَا طَائِلَةَ مِنْهَا} â†' \textarabic{مَا إِلهَا فَايْدَة} (maË Ê"ilha faËjde) — "it has no benefit" using possessive construction
\item \textarabic{طَائِلَةَ} â†' \textarabic{فَايْدَة} (faËjde) — more common dialectal word for "benefit"
\item \textarabic{فِي الدُّنْيَا} â†' \textarabic{بِالدُّنْيِي} (bi-d-diniË) — dialectal pronunciation with \textarabic{بِ} and final /iË/
\end{itemize}

\subsection{Alternative Dialectal Versions}

\textbf{Egyptian:}\\
\textarabic{لِيهْ بِيْحَشُّوا فِي دِمَاغِي لُغَة مَالهَاشْ لَازْمَة فِي الدُّنْيَا؟}\\
{[liË bijaħaʃʃuË fiË dimaËÉ£iË luÉ£a malhaËʃ laËzma fi-d-dinjÄ]}

\textbf{Gulf (Khaleeji):}\\
\textarabic{لِيشْ يِحْشِرُون بِرَاسِي لُغَة مَا لهَا فَايْدَة بِالدُّنْيَا؟}\\
{[liËʃ jiħʃiruËn biraËsiË luÉ£a maË lahaË faËjda bi-d-dinjÄ]}

% ======================== Additional Learning Notes ========================
\section{Additional Learning Notes}

\begin{tcolorbox}[colback=boxcolor,colframe=accentcolor,title=\textbf{Cultural and Pragmatic Context},breakable]

\textbf{Educational Critique:} This phrase reflects a common student frustration across Arabic-speaking cultures regarding mandatory curriculum subjects perceived as impractical or irrelevant to real life.

\textbf{Rhetorical Force:} The use of \textarabic{لِمَ} (why) rather than \textarabic{لِمَاذَا} creates a more literary, serious tone, elevating the complaint beyond casual grumbling to principled objection.

\textbf{Passive Voice Strategy:} By using the passive \textarabic{أُقْحَمُ}, the speaker emphasizes victimhood and lack of agency, implying criticism of an educational system or authority.

\textbf{Metaphor Analysis:} \textarabic{فِي مُخِّي} (in my brain) is more visceral than saying "taught to me" — it emphasizes forced insertion of unwanted knowledge.

\textbf{Universal vs. Absolute:} The phrase \textarabic{فِي الدُّنْيَا} (in the world) creates hyperbolic emphasis — not just "no benefit to me" but "no benefit anywhere in existence."

\textbf{Modern Relevance:} This type of complaint has become increasingly common in discussions about education reform, particularly regarding classical language instruction or theoretical subjects.
\end{tcolorbox}

\subsection{Memory Tips}
\begin{enumerate}
\item \textbf{Question Word:} \textit{\textarabic{لِمَ} = \textarabic{لِ} (for) + \textarabic{مَا} (what) = "for what?" = why?}

\item \textbf{Passive Recognition:} Form IV passive pattern: \textit{\textarabic{أُفْعَلُ} (Ê"ufÊ•alu)} — notice the \textarabic{أُ} (u) prefix and \textarabic{َ} (a) on middle radical

\item \textbf{Categorical Negation:} Remember \textit{\textarabic{لَا} + indefinite accusative noun = absolute negation}
   \begin{itemize}
   \item \textarabic{لَا رَجُلَ} — no man (at all)
   \item \textarabic{لَا فَائِدَةَ} — no benefit (whatsoever)
   \end{itemize}

\item \textbf{Body Part + Learning:} Common metaphor pattern in Arabic:
   \begin{itemize}
   \item \textarabic{فِي مُخِّي} — in my brain (forced memorization)
   \item \textarabic{فِي رَأْسِي} — in my head (thoughts)
   \item \textarabic{عَلَى لِسَانِي} — on my tongue (about to say)
   \end{itemize}

\item \textbf{Complaint Structure:} \textit{\textarabic{لِمَ} + passive verb + negative consequence} is a productive pattern for expressing frustration

\item \textbf{\textarabic{طَائِلَة} Recognition:} This formal word appears mainly in the fixed expression \textarabic{لَا طَائِلَةَ مِنْ} (no benefit from)
\end{enumerate}

\subsection{Related Vocabulary Family}

\begin{multicols}{2}
\textbf{Learning/Education:}\\
\textarabic{تَعَلُّم} — learning [taÊ•allum]\\
\textarabic{دِرَاسَة} — study [diraËsa]\\
\textarabic{حِفْظ} — memorization [ħifðˤ]\\
\textarabic{تَعْلِيم} — teaching [taÊ•liËm]\\
\textarabic{مُقَرَّر} — curriculum [muqarrar]\\
\textarabic{مَادَّة} — subject [mÄdda]\\
\textarabic{طَالِب} — student [tˤaËlib]\\[0.5em]

\textbf{Coercion/Force:}\\
\textarabic{إِجْبَار} — coercion [Ê"idÊ'baËr]\\
\textarabic{إِرْغَام} — compulsion [Ê"irÉ£aËm]\\
\textarabic{إِلْزَام} — obligation [Ê"ilzaËm]\\
\textarabic{فَرْض} — imposition [fardˤ]\\
\textarabic{اضْطِرَار} — necessity [idˤtˤiraËr]\\[0.5em]

\textbf{Benefit/Utility:}\\
\textarabic{فَائِدَة} — benefit [faËÊ"ida]\\
\textarabic{نَفْع} — utility [nafÊ•]\\
\textarabic{جَدْوَى} — usefulness [dÊ'adwaË]\\
\textarabic{قِيمَة} — value [qiËma]\\
\textarabic{أَهَمِّيَّة} — importance [Ê"ahammijja]\\
\textarabic{مَنْفَعَة} — advantage [manfaÊ•a]\\[0.5em]

\textbf{Futility/Uselessness:}\\
\textarabic{عَبَث} — futility [Ê•abaθ]\\
\textarabic{بَاطِل} — vain/false [baËtˤil]\\
\textarabic{لَا جَدْوَى} — useless [laË dÊ'adwaË]\\
\textarabic{عَقِيم} — sterile/fruitless [Ê•aqiËm]\\
\textarabic{بِلَا طَائِل} — pointless [bilaË tˤaËÊ"il]\\[0.5em]

\textbf{Brain/Mind:}\\
\textarabic{مُخّ} — brain [muxx]\\
\textarabic{دِمَاغ} — brain [dimaËÉ£]\\
\textarabic{عَقْل} — mind/intellect [Ê•aql]\\
\textarabic{ذِهْن} — mind [ðihn]\\
\textarabic{فِكْر} — thought [fikr]\\
\textarabic{ذَاكِرَة} — memory [ðaËkira]\\[0.5em]

\textbf{Language:}\\
\textarabic{لُغَة} — language [luÉ£a]\\
\textarabic{لِسَان} — tongue/language [lisaËn]\\
\textarabic{لَهْجَة} — dialect [lahdÊ'a]\\
\textarabic{كَلَام} — speech [kalaËm]\\
\textarabic{لُغَة أَجْنَبِيَّة} — foreign language\\
\textarabic{لُغَة أُمّ} — mother tongue
\end{multicols}

\subsection{Grammar Drills}

\subsubsection*{Practice: Form IV Passive Conjugation}
Transform these Form IV active verbs to passive, 1st person singular:

\begin{enumerate}
\item \textarabic{أَرْسَلَ} (to send) â†' \underline{\hspace{3cm}} [Ê"ursalu — I am sent]
\item \textarabic{أَكْرَمَ} (to honor) â†' \underline{\hspace{3cm}} [Ê"ukramu — I am honored]
\item \textarabic{أَخْرَجَ} (to expel) â†' \underline{\hspace{3cm}} [Ê"uxradÊ'u — I am expelled]
\item \textarabic{أَعْلَنَ} (to announce) â†' \underline{\hspace{3cm}} [Ê"uÊ•lanu — I am announced]
\end{enumerate}

\subsubsection*{Practice: \textarabic{لَا} Negation}
Complete these phrases with appropriate accusative forms:

\begin{enumerate}
\item \textarabic{لَا ـــــــــ فِي الْبَيْتِ} (no man) [\underline{رَجُلَ}]
\item \textarabic{لَا ـــــــــ لَهُ} (no friend) [\underline{صَدِيقَ}]
\item \textarabic{لَا ـــــــــ مِنْ هَذَا} (no benefit) [\underline{فَائِدَةَ}]
\item \textarabic{لَا ـــــــــ فِي ذَلِكَ} (no doubt) [\underline{شَكَّ}]
\end{enumerate}

\subsubsection*{Practice: Building Rhetorical Questions}
Create complaint questions using the pattern: \textit{\textarabic{لِمَ} + passive verb + negative clause}

\textit{Example:} \textarabic{لِمَ أُعَاقَبُ عَلَى خَطَإٍ لَمْ أَفْعَلْهُ؟}\\
(Why am I punished for an error I didn't commit?)

\section{Conclusion}

This phrase exemplifies sophisticated Arabic rhetorical complaint structures, combining:
\begin{itemize}
\item Passive voice to emphasize lack of agency
\item Categorical negation for absolute statements
\item Vivid metaphor for cognitive imposition
\item Formal register to lend gravity to complaint
\item Hyperbolic scope (\textarabic{فِي الدُّنْيَا}) for emphasis
\end{itemize}

The structure is highly productive and can be adapted to critique any mandatory activity perceived as useless — a pattern particularly relevant in educational and bureaucratic contexts throughout the Arabic-speaking world.

\end{document}